% esl-architecture-block — MPSoC-style two-tile block/bus architecture: each tile pairs a CPU core, an accelerator core, and a config-register block, joined by an intra-tile bus, with the two tiles joined by a dashed inter-tile NoC bus through a shared I/O hub.
% THESIS: Local (intra-tile) and global (inter-tile/NoC) interconnect are visually and structurally distinct, so per-tile compute reads separately from cross-tile communication.
% Concept grafted from: tsarilp-pub doc/common/tikzsetup.tex conventions (grafted, not a fork)
% Intent record: intent.yaml (contract §4 shape; ADR-0005 D3 §4 row).
\documentclass[border=8pt]{standalone}

% --- packages (mirror these in template.meta.json "requires") ---
\usepackage[dvipsnames]{xcolor}
\usepackage{tikz}
\usetikzlibrary{calc, shapes, positioning, backgrounds, fit, patterns, arrows.meta}

% --- palette: ESL domain-semantic tokens (contract §2, ADR-0005 D7 -- extends
% opentikz's ot* palette, does not remap onto it; see
% reference/color-palettes/color-palettes.md "ESL domain-semantic tokens").
% Originally grafted verbatim from tsarilp's tikzsetup.tex color names
% (Aquamarine/NavyBlue/ForestGreen/Orange); now bound through the named
% palette tokens tools/validate.py enforces instead of the raw dvipsnames.
\colorlet{core-accent}{Aquamarine}       % processing-core fill accent
\colorlet{process-accent}{NavyBlue}      % flow/transform-stage (accelerator) accent
\colorlet{param-accent}{ForestGreen}     % parameter/artifact box accent
\colorlet{intra-bus}{Orange}             % intra-tile interconnect
\colorlet{inter-bus}{NavyBlue}           % inter-tile interconnect
\colorlet{digitaldomain}{gray}           % digital words/logic -- shared I/O terminal

\begin{document}
% ==== parameters (edit these) ============================================
\def\tilesep{2.6}     % horizontal gap between tile boundaries (cm)
\def\coresize{50pt}   % core box minimum height
\def\corewidth{64pt}  % core box minimum width
% labels
\def\tileonelabel{Tile 0}
\def\tiletwolabel{Tile 1}
\def\coreonelabel{Core}
\def\coretwolabel{Core}
\def\acconelabel{Accel}
\def\acctwolabel{Accel}
\def\cfgonelabel{Cfg\\Regs}
\def\cfgtwolabel{Cfg\\Regs}
\def\iolabel{I/O}
\def\interbuslabel{Inter-tile bus (NoC)}
% =========================================================================

\begin{tikzpicture}[
    node distance=8pt and \tilesep cm,
    >=Stealth,
    % --- grafted styles (tikzsetup.tex naming/coloring kept verbatim) ---
    PE/.style={draw, thick, minimum size=35pt, align=center},
    core/.style={PE, draw=core-accent, fill=core-accent!20,
      minimum width=\corewidth, minimum height=\coresize},
    process/.style={draw=process-accent, fill=process-accent!10, very thick,
      minimum width=\corewidth, minimum height=0.7*\coresize, align=center,
      font=\scriptsize},
    parameters/.style={chamfered rectangle,
      chamfered rectangle corners={north east},
      draw=param-accent, fill=param-accent!20, very thick,
      minimum width=0.85*\corewidth, minimum height=0.55*\coresize,
      align=center, font=\scriptsize},
    terminal node/.style={draw=digitaldomain, fill=digitaldomain!15, text=digitaldomain!70!black, thick,
      circle, minimum size=22pt, font=\scriptsize},
    intrabus/.style={draw, intra-bus, thick},
    interbus/.style={draw, inter-bus, thick, dashed},
    % --- fit-margins helper, grafted verbatim from tikzsetup.tex ---
    fit margins/.style={/tikz/afit/.cd,#1,
      /tikz/.cd,
      inner xsep=\pgfkeysvalueof{/tikz/afit/left}+\pgfkeysvalueof{/tikz/afit/right},
      inner ysep=\pgfkeysvalueof{/tikz/afit/top}+\pgfkeysvalueof{/tikz/afit/bottom},
      xshift=-\pgfkeysvalueof{/tikz/afit/left}+\pgfkeysvalueof{/tikz/afit/right},
      yshift=-\pgfkeysvalueof{/tikz/afit/bottom}+\pgfkeysvalueof{/tikz/afit/top}},
    afit/.cd, left/.initial=6pt, right/.initial=6pt,
    bottom/.initial=6pt, top/.initial=14pt,
  ]

  % ---- Tile 0 ----
  \node[core] (core1) {\coreonelabel};
  \node[process, below=of core1] (acc1) {\acconelabel};
  \node[parameters, right=10pt of core1] (cfg1) {\cfgonelabel};
  \draw[intrabus] (core1) -- (acc1);
  \draw[intrabus] (core1) -- (cfg1);

  \begin{scope}[on background layer]
    \node[draw=core-accent!70!black, dashed, rounded corners=4pt,
      fit margins={left=8pt,right=8pt,top=16pt,bottom=8pt},
      fit=(core1)(acc1)(cfg1)] (tile1) {};
    \node[font=\scriptsize\bfseries, text=core-accent!60!black,
      anchor=north west] at ([xshift=2pt,yshift=-2pt]tile1.north west)
      {\tileonelabel};
  \end{scope}

  % ---- Tile 1 ----
  \node[core, right=\tilesep cm of tile1] (core2) {\coretwolabel};
  \node[process, below=of core2] (acc2) {\acctwolabel};
  \node[parameters, right=10pt of core2] (cfg2) {\cfgtwolabel};
  \draw[intrabus] (core2) -- (acc2);
  \draw[intrabus] (core2) -- (cfg2);

  \begin{scope}[on background layer]
    \node[draw=core-accent!70!black, dashed, rounded corners=4pt,
      fit margins={left=8pt,right=8pt,top=16pt,bottom=8pt},
      fit=(core2)(acc2)(cfg2)] (tile2) {};
    \node[font=\scriptsize\bfseries, text=core-accent!60!black,
      anchor=north west] at ([xshift=2pt,yshift=-2pt]tile2.north west)
      {\tiletwolabel};
  \end{scope}

  % ---- shared I/O terminal + inter-tile bus ----
  % calc midpoint idiom (tikzsetup.tex convention: ($(a)!0.5!(b)$)), bound to a
  % named \coordinate so it can be a "right=of"/"above=of" positioning anchor.
  \coordinate (tilemid) at ($(tile1.north)!0.5!(tile2.north)$);
  \node[terminal node, above=18pt of tilemid] (io) {\iolabel};
  \draw[interbus] (tile1.north) -- (io);
  \draw[interbus] (tile2.north) -- (io);
  \node[font=\tiny, text=inter-bus!70!black, above=6pt of io] {\interbuslabel};

\end{tikzpicture}
\end{document}
